\medskip

Dans l'espace muni d'un repère orthonormé \Oijk, on considère les points 
\begin{center}A(0~;~4~;~16), \quad B$(0~;~4~;~- 10)$, \quad C$(4~;~-8~;~0)$\quad et\quad  K(0~;~4~;~3).\end{center}

On définit la sphère $S$ de centre K et de rayon 13 comme l'ensemble des points M tels que KM $= 13$.

\medskip

\begin{enumerate}
\item 
	\begin{enumerate}
		\item Vérifier que le point C appartient à la sphère $S$
		\item Montrer que le triangle ABC est rectangle en C.
	\end{enumerate}
\item 
	\begin{enumerate}
		\item Montrer que le vecteur $\vect{n}\begin{pmatrix}3\\1\\0\end{pmatrix}$ est un vecteur normal au plan (ABC).
		\item Déterminer une équation cartésienne du plan (ABC).
	\end{enumerate}
\item On admet que la sphère $S$ coupe l'axe des abscisses en deux points, l'un ayant une abscisse positive et l'autre une abscisse négative.

On note D celui qui a une abscisse positive.
	\begin{enumerate}
		\item Montrer que le point D a pour coordonnées (12~;~0~;~0).
		\item Donner une représentation paramétrique de la droite $\Delta$ passant par D et perpendiculaire au plan (ABC).
		\item  Déterminer la distance du point D au plan (ABC).
	\end{enumerate}
\item Calculer une valeur approchée, à l'unité de volume près, du volume du tétraèdre ABCD.

\emph{On rappelle la formule du volume V d'un tétraèdre}
\[V = \dfrac13 \times \mathcal{B} \times h.\]

\emph{où $\mathcal{B}$ est l'aire d'une base et $h$ la hauteur associée}.
\end{enumerate}

\bigskip
